\scanpage{32}
\mspage{13}

En résumé :

\uline{Proposition 2} : \unsure{Soient|Soit} $M$ un monoïde, \deleted{On suppose} $\varphi$ et $\psi$ deux éléments de $M$.

On suppose :

a) $\psi$ est un élément minimal de $M$, i.e pour \deleted{tout} $u\in M$, existe $\unsure{\Theta|\theta}(u)\in M$ avec $\psi=u\Theta(u)$. (Donc $M$ i.e. $Y=B_M$ est pseudo-cofiltrant, i.e. $Y^\wedge$ est loc. irréductible.)

b) $\psi\sim1$ pour la relation d'équivalence invariante à droite $R_\varphi$ définie par $\varphi$
\[
(u\overset{R_\varphi}{\sim}v\quad\Longleftrightarrow\quad\exists\,p,q\in\mathbb N\ \text{avec}\ \varphi^p u=\varphi^q v),
\]
i.e $\exists\,\alpha,\beta\in\mathbb N$ avec $\varphi^\alpha\psi=\varphi^\beta\deleted{\unsure{\psi}}$.

c) Si $\deleted{p,q}\unsure{r,s|p,q}\in\mathbb N$ sont tels que $\varphi^{\deleted{\unsure{p}}r}\psi=\varphi^s\psi$, alors $r=s$.

Soit $Y=B_M$, $F=M/\unsure{\mathcal R_\varphi|R_\varphi}$ considéré comme $M$-ensemble à droite i.e. comme objet de $Y^\wedge$, $X=Y_{/F}$. Alors l'objet $e=\dot 1$ de $X$ est à la fois un plus \deleted{petit}\inserted{grand} et un plus petit élément de $X$.

De plus, $\deleted{\pi_{1}(X,e)}$ il existe un hom surjectif $\pi_1(X,e)\longrightarrow\mathbb Z$, en particulier $X$ n'est pas $1$-connexe.

% archive sheet number: 31
